\chapter{Depth-First Search, Timestamps, and Structure}

\chapterline

Breadth-first search fans out in rings: it touches everything one step away, then everything two steps away, and the geometry it discovers is the geometry of shortest hops. Depth-first search does the opposite. It commits. It walks as far down one branch as it can, and only when the branch dead-ends does it back up to the most recent fork and try the next option. That stubbornness looks like a worse search strategy --- and for shortest paths it is --- but the recursion stack of a DFS is a remarkably sensitive instrument. The \emph{times} at which DFS first enters and finally leaves each vertex encode a nesting structure, and that nesting is the hidden skeleton behind topological sort, strongly connected components, bridges, articulation points, biconnectivity, Euler tours, and the dominator tree. Almost every ``clever linear-time graph trick'' in the later chapters is, underneath, a fact about DFS timestamps.

This chapter builds that instrument carefully. We run DFS by hand on a small graph, record the discovery and finish times, prove the \emph{parenthesis theorem} that governs how the time intervals nest, classify every edge of the graph as tree, back, forward, or cross, and use that classification to detect cycles in both directed and undirected graphs. Throughout, our goal is to show why the discovery and finishing timestamps form a rigid algebraic structure, and how later algorithms translate that structure directly into cycle certificates, topological orderings, and components.

\begin{focusbox}
By the end you should be able to: run DFS and annotate a graph with $d[\cdot]$ and $f[\cdot]$; state and prove the parenthesis and white-path theorems; classify an edge from the timestamps of its endpoints alone; explain why a back edge means a cycle in a digraph but why undirected graphs need the parent check; and name three later algorithms that are ``just'' a reading of these timestamps.
\end{focusbox}

\section{The walk: intuition and a hand trace}
\label{ch03:intuition}

\begin{intuitionbox}
Imagine exploring a cave system with a ball of string and a lamp. From the entrance you pick a passage and walk down it, unspooling string. At every junction you pick an unexplored passage and keep going. When you hit a chamber with no new exits, you reel the string back to the last junction that still has an unexplored passage and dive down that one. The string is the recursion stack. The moment you \emph{first} enter a chamber is its discovery time; the moment you reel past it for the last time, having exhausted everything reachable beyond it, is its finish time. DFS is exactly this: go deep, backtrack minimally, and timestamp both events.
\end{intuitionbox}

We maintain a global clock that starts at $0$ and ticks by one each time \emph{anything} happens --- each discovery and each finish. Every vertex $u$ gets two stamps:
\begin{itemize}
  \item $d[u]$, the clock value when $u$ is first reached (colored gray, ``on the stack'');
  \item $f[u]$, the clock value when DFS finishes exploring everything reachable from $u$ and pops it (colored black, ``done'').
\end{itemize}
Vertices start white (undiscovered). With $n$ vertices, the clock runs from $1$ to $2n$, and every integer in that range is used exactly once: each vertex contributes one discovery tick and one finish tick.

\subsection{A directed graph to trace}

Consider the directed graph $\Gph$ on six vertices $\{1,2,3,4,5,6\}$ with edges
\[
1\to2,\quad 1\to4,\quad 2\to3,\quad 3\to1,\quad 4\to3,\quad 5\to4,\quad 5\to6,\quad 6\to6.
\]

\begin{figure}[H]
\centering
\begin{tikzpicture}[
  >={Stealth[round]}, node distance=20mm,
  vtx/.style={circle,draw=nodeblue,very thick,fill=blue!4,minimum size=9mm,font=\bfseries},
  every edge/.style={draw=edgegray,thick}]
  \node[vtx] (1) {1};
  \node[vtx] (2) [right=of 1] {2};
  \node[vtx] (3) [below=of 2] {3};
  \node[vtx] (4) [below=of 1] {4};
  \node[vtx] (5) [left=22mm of 1] {5};
  \node[vtx] (6) [below=of 5] {6};
  \path[->]
    (1) edge (2)
    (1) edge (4)
    (2) edge (3)
    (3) edge[bend left=12] (1)
    (4) edge (3)
    (5) edge (4)
    (5) edge (6)
    (6) edge[loop below,looseness=8] (6);
\end{tikzpicture}
\caption{The running directed example. Adjacency lists are kept in increasing vertex order, and the outer loop visits roots $1,2,\dots,6$ in order.}
\label{ch03:fig-digraph}
\end{figure}

We always break ties the same way so the trace is reproducible: the outer loop considers vertices $1,2,\dots,6$ as potential roots, and each adjacency list is scanned in increasing order. Start the clock at $0$.

\begin{labbox}
Follow along with pencil. Maintain a stack (top on the right) and a clock. Each line below is one clock tick.
\begin{enumerate}\itemsep0.15em
  \item Root $1$ is white. Discover $1$: $d[1]=1$. Stack: $[1]$.
  \item $1$'s first out-neighbor is $2$ (white). Discover $2$: $d[2]=2$. Stack: $[1,2]$.
  \item $2$'s out-neighbor is $3$ (white). Discover $3$: $d[3]=3$. Stack: $[1,2,3]$.
  \item $3$'s out-neighbor is $1$, which is gray (on the stack). This is a \emph{back edge}; do not recurse. $3$ has no other out-edges. Finish $3$: $f[3]=4$. Pop. Stack: $[1,2]$.
  \item Back at $2$: no more out-edges. Finish $2$: $f[2]=5$. Pop. Stack: $[1]$.
  \item Back at $1$: next out-neighbor is $4$ (white). Discover $4$: $d[4]=6$. Stack: $[1,4]$.
  \item $4$'s out-neighbor is $3$, which is \emph{black} (finished). Edge $4\to3$ is a \emph{cross} edge (we will justify the name shortly). $4$ has no other out-edges. Finish $4$: $f[4]=7$. Pop. Stack: $[1]$.
  \item Back at $1$: no more out-edges. Finish $1$: $f[1]=8$. Pop. Stack empty. Tree rooted at $1$ is done.
  \item Outer loop advances. Vertices $2,3,4$ are already black; skip. Vertex $5$ is white. Discover $5$: $d[5]=9$. Stack: $[5]$.
  \item $5$'s first out-neighbor is $4$ (black). Edge $5\to4$ is a cross edge. Next out-neighbor $6$ (white). Discover $6$: $d[6]=10$. Stack: $[5,6]$.
  \item $6$'s only out-edge is the self-loop $6\to6$ ($6$ is gray): a back edge. Finish $6$: $f[6]=11$. Pop. Stack: $[5]$.
  \item Finish $5$: $f[5]=12$. Pop. Outer loop finds $6$ black, ends. Clock at $12=2\cdot 6$. Done.
\end{enumerate}
\end{labbox}

Collecting the stamps as $(d,f)$ pairs:
\[
1{:}(1,8)\quad 2{:}(2,5)\quad 3{:}(3,4)\quad 4{:}(6,7)\quad 5{:}(9,12)\quad 6{:}(10,11).
\]

\begin{figure}[H]
\centering
\begin{tikzpicture}[
  >={Stealth[round]}, node distance=22mm,
  vtx/.style={circle,draw=nodeblue,very thick,fill=blue!4,minimum size=11mm,font=\small\bfseries},
  tree/.style={draw=structurepurple,very thick,->},
  back/.style={draw=trapred,thick,->,dashed},
  cross/.style={draw=targetgreen,thick,->,dotted},
  stamp/.style={font=\scriptsize,text=black!70}]
  \node[vtx,label={[stamp]above:{1/8}}] (1) {1};
  \node[vtx,label={[stamp]above:{2/5}}] (2) [right=of 1] {2};
  \node[vtx,label={[stamp]right:{3/4}}] (3) [below=of 2] {3};
  \node[vtx,label={[stamp]left:{6/7}}] (4) [below=of 1] {4};
  \node[vtx,label={[stamp]above:{9/12}}] (5) [left=26mm of 1] {5};
  \node[vtx,label={[stamp]left:{10/11}}] (6) [below=of 5] {6};
  \path
    (1) edge[tree] (2)
    (2) edge[tree] (3)
    (1) edge[tree] (4)
    (5) edge[tree] (6)
    (3) edge[back,bend left=12] (1)
    (6) edge[back,loop below,looseness=8] (6)
    (4) edge[cross] (3)
    (5) edge[cross] (4);
\end{tikzpicture}
\caption{The DFS forest with $d/f$ stamps. Tree edges (purple, solid), back edges (red, dashed), cross edges (green, dotted). Two trees: one rooted at $1$, one at $5$.}
\label{ch03:fig-forest}
\end{figure}

Notice the two roots $1$ and $5$ produced two trees; together they form the \emph{DFS forest}. Notice too that the dashed back edges $3\to1$ and $6\to6$ each close a cycle, and that no cross edge ever points ``into'' a still-open subtree. Those are not accidents of this example --- they are theorems, and we will prove them.

\section{The algorithm}
\label{ch03:algorithm}

Here is DFS with timestamps and edge classification, written recursively. The wrapper loops over all vertices so that disconnected graphs (and digraphs not reachable from one root) are fully covered, yielding a forest rather than a single tree.

\begin{algorithm}[H]
\DontPrintSemicolon
\SetKwInOut{Input}{input}\SetKwInOut{Output}{output}
\Input{Graph $\Gph=(V,E)$ as adjacency lists; directed or undirected.}
\Output{Stamps $d[\cdot],f[\cdot]$, parent $\pi[\cdot]$, and an edge-type label for each edge.}
\BlankLine
\ForEach{$u\in V$}{$\mathit{color}[u]\gets\textsc{white}$; $\pi[u]\gets\textsc{nil}$\;}
$\mathit{time}\gets 0$\;
\ForEach{$u\in V$ \textnormal{in fixed order}}{
  \lIf{$\mathit{color}[u]=\textsc{white}$}{$\textsc{DFS-Visit}(u)$}
}
\BlankLine
\SetKwProg{Fn}{procedure}{}{}
\Fn{\textsc{DFS-Visit}$(u)$}{
  $\mathit{time}\gets\mathit{time}+1$;\quad $d[u]\gets\mathit{time}$;\quad $\mathit{color}[u]\gets\textsc{gray}$\;
  \ForEach{$v\in\mathit{Adj}[u]$}{
    \uIf{$\mathit{color}[v]=\textsc{white}$}{
      classify $(u,v)$ as \textbf{tree}\;
      $\pi[v]\gets u$;\quad $\textsc{DFS-Visit}(v)$\;
    }
    \uElseIf{$\mathit{color}[v]=\textsc{gray}$}{
      classify $(u,v)$ as \textbf{back}\;
    }
    \Else{
      \lIf{$d[u]<d[v]$}{classify $(u,v)$ as \textbf{forward}}
      \lElse{classify $(u,v)$ as \textbf{cross}}
    }
  }
  $\mathit{time}\gets\mathit{time}+1$;\quad $f[u]\gets\mathit{time}$;\quad $\mathit{color}[u]\gets\textsc{black}$\;
}
\caption{Depth-first search with timestamps and edge classification}
\label{ch03:alg-dfs}
\end{algorithm}

The colors are the whole accounting trick. \textsc{white} means undiscovered, \textsc{gray} means discovered-but-not-finished (equivalently, \emph{currently on the recursion stack}), and \textsc{black} means finished. An edge $(u,v)$ encountered while scanning $u$'s list is classified purely by the color of $v$ at that instant, with one numeric tiebreak between forward and cross. We prove that this classification is correct and complete in Section~\ref{ch03:edge-types}.

\begin{trapbox}
On large graphs the recursion in Algorithm~\ref{ch03:alg-dfs} can blow the call stack: a path graph on $n=10^6$ vertices recurses a million deep, and most language runtimes overflow well before that. In production code DFS is run with an explicit stack. The subtlety is that a correct iterative DFS must record \emph{where it was} in each adjacency list (an iterator per stack frame) and must run the ``finish'' bookkeeping when a frame is popped --- not when it is first pushed. Iterative versions that timestamp on push compute $d[\cdot]$ correctly but get $f[\cdot]$ and the parenthesis structure wrong. We give a correct iterative skeleton in Exercise~\ref{ch03:ex-iterative}.
\end{trapbox}

\subsection{The iterative skeleton}
\label{ch03:iterative}

Because the recursion above is the part most likely to break on real inputs, it is worth seeing the explicit-stack version in full once, rather than waving at it. The trick that makes it match the recursive timestamps exactly is to store, with each stacked vertex, an \emph{index} into its adjacency list --- the cursor marking how far we have scanned --- and to emit the finish stamp only when a frame is finally popped because its cursor has run off the end of the list.

\begin{algorithm}[H]
\DontPrintSemicolon
\SetKwInOut{Input}{input}\SetKwInOut{Output}{output}
\Input{Graph $\Gph=(V,E)$ as adjacency lists; a start vertex $s$.}
\Output{Stamps $d[\cdot],f[\cdot]$ identical to the recursive DFS rooted at $s$.}
\BlankLine
$\mathit{time}\gets 0$\;
push frame $(s,\,0)$ onto stack $S$\;
$\mathit{time}\gets\mathit{time}+1$;\quad $d[s]\gets\mathit{time}$;\quad $\mathit{color}[s]\gets\textsc{gray}$\;
\While{$S$ \textnormal{nonempty}}{
  let $(u,i)$ be the top frame of $S$\;
  \uIf{$i<\abs{\mathit{Adj}[u]}$}{
    $v\gets\mathit{Adj}[u][i]$\;
    increment the stored cursor of the top frame to $i+1$\;
    \uIf{$\mathit{color}[v]=\textsc{white}$}{
      $\pi[v]\gets u$\;
      $\mathit{time}\gets\mathit{time}+1$;\quad $d[v]\gets\mathit{time}$;\quad $\mathit{color}[v]\gets\textsc{gray}$\;
      push frame $(v,\,0)$ onto $S$\;
    }
  }
  \Else{
    $\mathit{time}\gets\mathit{time}+1$;\quad $f[u]\gets\mathit{time}$;\quad $\mathit{color}[u]\gets\textsc{black}$\;
    pop $S$\;
  }
}
\caption{Iterative DFS with a per-frame adjacency cursor}
\label{ch03:alg-iterative}
\end{algorithm}

The single most important line is the \textbf{else} branch: a vertex is finished and blackened only when its cursor has walked off the end of its list, i.e.\ when control would have returned from the recursive call. Push the start onto the stack, advance the top frame's cursor one neighbor at a time, and you reproduce the recursion's control flow exactly --- including the order in which finishes happen, which is the whole content of the parenthesis structure. Run this skeleton on Figure~\ref{ch03:fig-digraph} starting at $1$ and the stamps come out $1{:}(1,8),\,2{:}(2,5),\,3{:}(3,4),\,4{:}(6,7)$, matching the recursive trace tick for tick. Edge classification slots in at the point where $v$ is examined, using the same color test as Algorithm~\ref{ch03:alg-dfs}.

\subsection{Complexity}

\begin{theorem}[Linear time]\label{ch03:thm-linear}
On a graph with $n=\abs{V}$ vertices and $m=\abs{E}$ edges given as adjacency lists, Algorithm~\ref{ch03:alg-dfs} runs in $\Th(n+m)$ time and $\Th(n)$ working space beyond the input.
\end{theorem}

\begin{proof}
\textsc{DFS-Visit} is called exactly once per vertex: it is called only on \textsc{white} vertices, and its first action is to color the vertex \textsc{gray}, so no vertex is visited twice. The initialization loop and the outer loop are $\Th(n)$. Inside the single visit of $u$, the \textbf{for} loop scans $\mathit{Adj}[u]$ once, doing $\Oh(1)$ work per incident edge. Summing the list lengths over all vertices counts each directed edge once, and each undirected edge twice (once from each endpoint), so the total edge-scanning work is $\Th(m)$ either way. The two timestamp updates per vertex add $\Th(n)$. Total: $\Th(n+m)$. The color, $d$, $f$, and $\pi$ arrays are $\Th(n)$; the recursion stack is at most $n$ deep, also $\Th(n)$.
\end{proof}

\begin{complexitybox}
\textbf{DFS (adjacency lists).} Time $\Th(n+m)$. Extra space $\Th(n)$ (arrays plus stack depth $\le n$). On an adjacency \emph{matrix} the edge scan becomes $\Th(n^2)$ regardless of $m$, since finding the neighbors of $u$ costs $\Th(n)$. The timestamps, parent pointers, and edge labels are all produced within this same bound --- the structure is free.
\end{complexitybox}

\section{The parenthesis theorem}
\label{ch03:parenthesis}

The single most useful fact about DFS is that the time intervals $[d[u],f[u]]$ behave like well-formed parentheses: any two are either disjoint or nested, never partially overlapping. Write the interval of $u$ as $I_u=[d[u],f[u]]$.

\begin{intuitionbox}
Think of $d[u]$ as printing an open paren tagged $u$ and $f[u]$ as printing the matching close paren tagged $u$. Because DFS is a stack discipline --- you cannot finish $u$ until everything you started after entering $u$ has finished --- the printed string is a balanced parenthesization. For our trace the string is
\[
\texttt{(}_1\,\texttt{(}_2\,\texttt{(}_3\,\texttt{)}_3\,\texttt{)}_2\,\texttt{(}_4\,\texttt{)}_4\,\texttt{)}_1\,\texttt{(}_5\,\texttt{(}_6\,\texttt{)}_6\,\texttt{)}_5
\]
and you can read the entire tree shape off the nesting: $3$ is inside $2$ is inside $1$; $4$ is inside $1$ but beside $2$; $6$ is inside $5$.
\end{intuitionbox}

\begin{center}
\begin{tikzpicture}[
    >=Stealth,
    ruler/.style={font=\scriptsize\ttfamily, text=gray!80!black},
    tick/.style={gray!35, thick},
    bar/.style={rounded corners=2pt, thick},
    labelcol/.style={font=\footnotesize\bfseries, text=gray!90!black, anchor=east}
]
  % Title
  \node[font=\footnotesize\scshape, text=gray!80!black, anchor=west] at (-0.5, 3.8) {Timeline of discovery and finishing intervals $I_u = [d[u], f[u]]$:};

  % Clock ruler
  \draw[thick, gray!60] (0.8, 3.2) -- (10.0, 3.2);
  \foreach \t/\x in {1/0.8, 2/1.6, 3/2.4, 4/3.2, 5/4.0, 6/4.8, 7/5.6, 8/6.4, 9/7.2, 10/8.0, 11/8.8, 12/9.6} {
    \draw[tick] (\x, 3.2) -- (\x, -1.7);
    \node[ruler, above=1pt] at (\x, 3.2) {\t};
  }

  % Bars for each vertex
  \node[labelcol] at (0.5, 2.6) {$u=1$:};
  \draw[bar, fill=sourceblue!15, draw=sourceblue!80] (0.8, 2.4) rectangle (6.4, 2.85);
  \node[font=\tiny\bfseries, text=sourceblue!95!black] at (3.6, 2.62) {$I_1 = [1, 8]$};

  \node[labelcol] at (0.5, 2.0) {$u=2$:};
  \draw[bar, fill=targetgreen!15, draw=targetgreen!80] (1.6, 1.8) rectangle (4.0, 2.25);
  \node[font=\tiny\bfseries, text=targetgreen!95!black] at (2.8, 2.02) {$I_2 = [2, 5]$};

  \node[labelcol] at (0.5, 1.4) {$u=3$:};
  \draw[bar, fill=quotientorange!15, draw=quotientorange!80] (2.4, 1.2) rectangle (3.2, 1.65);
  \node[font=\tiny\bfseries, text=quotientorange!95!black] at (2.8, 1.42) {$I_3 = [3,4]$};

  \node[labelcol] at (0.5, 0.8) {$u=4$:};
  \draw[bar, fill=structurepurple!15, draw=structurepurple!80] (4.8, 0.6) rectangle (5.6, 1.05);
  \node[font=\tiny\bfseries, text=structurepurple!95!black] at (5.2, 0.82) {$I_4 = [6,7]$};

  \node[labelcol] at (0.5, 0.2) {$u=5$:};
  \draw[bar, fill=sourceblue!15, draw=sourceblue!80] (7.2, 0.0) rectangle (9.6, 0.45);
  \node[font=\tiny\bfseries, text=sourceblue!95!black] at (8.4, 0.22) {$I_5 = [9, 12]$};

  \node[labelcol] at (0.5, -0.4) {$u=6$:};
  \draw[bar, fill=targetgreen!15, draw=targetgreen!80] (8.0, -0.6) rectangle (8.8, -0.15);
  \node[font=\tiny\bfseries, text=targetgreen!95!black] at (8.4, -0.38) {$I_6 = [10,11]$};

  % Parenthesis string track at the bottom
  \node[labelcol] at (0.5, -1.25) {parens:};
  \node[font=\scriptsize\ttfamily] at (0.8, -1.25) {$\mathbf{(}_1$};
  \node[font=\scriptsize\ttfamily] at (1.6, -1.25) {$\mathbf{(}_2$};
  \node[font=\scriptsize\ttfamily] at (2.4, -1.25) {$\mathbf{(}_3$};
  \node[font=\scriptsize\ttfamily] at (3.2, -1.25) {$\mathbf{)}_3$};
  \node[font=\scriptsize\ttfamily] at (4.0, -1.25) {$\mathbf{)}_2$};
  \node[font=\scriptsize\ttfamily] at (4.8, -1.25) {$\mathbf{(}_4$};
  \node[font=\scriptsize\ttfamily] at (5.6, -1.25) {$\mathbf{)}_4$};
  \node[font=\scriptsize\ttfamily] at (6.4, -1.25) {$\mathbf{)}_1$};
  \node[font=\scriptsize\ttfamily] at (7.2, -1.25) {$\mathbf{(}_5$};
  \node[font=\scriptsize\ttfamily] at (8.0, -1.25) {$\mathbf{(}_6$};
  \node[font=\scriptsize\ttfamily] at (8.8, -1.25) {$\mathbf{)}_6$};
  \node[font=\scriptsize\ttfamily] at (9.6, -1.25) {$\mathbf{)}_5$};
\end{tikzpicture}
\end{center}

\begin{theorem}[Parenthesis theorem]\label{ch03:thm-paren}
In any DFS of a (directed or undirected) graph, for every pair of vertices $u\ne v$ exactly one of the following holds:
\begin{enumerate}
  \item[(i)] $I_u$ and $I_v$ are disjoint, and neither of $u,v$ is a descendant of the other in the DFS forest;
  \item[(ii)] $I_u\subset I_v$ (strictly), and $u$ is a descendant of $v$;
  \item[(iii)] $I_v\subset I_u$ (strictly), and $v$ is a descendant of $u$.
\end{enumerate}
In particular the intervals are \emph{laminar}: no two of them properly cross.
\end{theorem}

\begin{proof}
All four numbers $d[u],f[u],d[v],f[v]$ are distinct (every clock value is unique). Without loss of generality assume $d[u]<d[v]$; the other case is symmetric by swapping names. There are two sub-cases depending on whether $d[v]<f[u]$.

\emph{Case A: $d[v]<f[u]$.} Then $v$ is discovered while $u$ is gray, i.e.\ while $u$ is still on the stack. By the stack discipline, $u$ cannot be popped (finished) until every vertex discovered after $u$ and still on the stack above it has been finished; $v$ is such a vertex. Hence $v$ finishes before $u$: $f[v]<f[u]$. Combining, $d[u]<d[v]<f[v]<f[u]$, so $I_v\subset I_u$. Moreover, $v$ was reached during the call $\textsc{DFS-Visit}(u)$, before that call returned, so by the recursive structure $v$ lies in the subtree rooted at $u$ --- $v$ is a descendant of $u$. This is case~(iii).

\emph{Case B: $d[v]>f[u]$.} Then $v$ is discovered only after $u$ has already finished, so the entire interval $I_u$ precedes $d[v]$: $d[u]<f[u]<d[v]<f[v]$. The intervals are disjoint. Neither is a descendant of the other: $v$ cannot be a descendant of $u$ because the call $\textsc{DFS-Visit}(u)$ had already returned before $v$ was touched; and $u$ cannot be a descendant of $v$ because $u$ was discovered first. This is case~(i).

Since exactly one of $d[v]<f[u]$ or $d[v]>f[u]$ holds (they are distinct), exactly one conclusion holds. Swapping the roles of $u,v$ when $d[v]<d[u]$ gives case~(ii).
\end{proof}

\begin{corollary}[Nesting of descendants]\label{ch03:cor-nest}
$v$ is a proper descendant of $u$ in the DFS forest if and only if $d[u]<d[v]<f[v]<f[u]$, i.e.\ if and only if $I_v\subset I_u$.
\end{corollary}

\begin{proof}
Immediate from Theorem~\ref{ch03:thm-paren}: containment of intervals and the ancestor/descendant relation coincide.
\end{proof}

This corollary is the workhorse. It turns the question ``is $v$ in $u$'s subtree?'' --- which sounds like it needs a tree traversal --- into a pair of integer comparisons on precomputed stamps. The Euler-tour technique of Chapter~21 is precisely this idea pushed to its limit, and the heavy-light and centroid decompositions lean on it too.

\subsection{The white-path theorem}

The parenthesis theorem describes the output; the white-path theorem describes \emph{why} a vertex ends up where it does, and it is the cleanest tool for proving later facts about reachability.

\begin{theorem}[White-path theorem]\label{ch03:thm-whitepath}
In a DFS forest of $\Gph$, vertex $v$ is a descendant of vertex $u$ if and only if, at the moment $d[u]$ when DFS first discovers $u$, there is a path from $u$ to $v$ consisting entirely of white (as-yet-undiscovered) vertices --- $u$ itself permitted as the white start.
\end{theorem}

\begin{proof}
($\Rightarrow$) Suppose $v$ is a descendant of $u$. Let $u=x_0,x_1,\dots,x_k=v$ be the tree path from $u$ down to $v$. Each $x_{i}$ is a tree-child of $x_{i-1}$, hence discovered during $x_{i-1}$'s visit, hence after $u$; so each $x_i$ with $i\ge 1$ is discovered strictly after time $d[u]$, meaning all of $x_1,\dots,x_k$ are white at time $d[u]$. Together with $u$ itself this is a white path from $u$ to $v$.

($\Leftarrow$) Suppose at time $d[u]$ there is a white path $u=x_0,x_1,\dots,x_k=v$. We show by induction on $i$ that each $x_i$ becomes a descendant of $u$ (with $u$ counted as a descendant of itself). The base case $x_0=u$ is trivial. Inductive step: assume $x_{i-1}$ is a descendant of $u$, so $d[u]\le d[x_{i-1}]$ and $I_{x_{i-1}}\subseteq I_u$. Consider the edge $x_{i-1}\to x_i$, and ask when $x_i$ is first discovered.

Suppose, for contradiction, $x_i$ is \emph{not} a descendant of $u$. By the parenthesis theorem the intervals $I_u$ and $I_{x_i}$ are then disjoint, and since $x_i$ is white at $d[u]$ it is discovered after $d[u]$; disjointness forces $I_u$ to come entirely before $I_{x_i}$, so $d[x_i]>f[u]\ge f[x_{i-1}]$. But the edge $x_{i-1}\to x_i$ is scanned during the visit of $x_{i-1}$, at a time strictly before $f[x_{i-1}]<d[x_i]$; at that scan instant $x_i$ has not yet been discovered, so it is white, and DFS would immediately discover $x_i$ as a tree-child of $x_{i-1}$. That makes $x_i$ a descendant of $x_{i-1}$, hence of $u$, contradicting our assumption. Therefore $x_i$ is a descendant of $u$, completing the induction. Taking $i=k$ gives that $v$ is a descendant of $u$.
\end{proof}

\begin{whybox}
The white-path theorem says the DFS subtree of $u$ is \emph{exactly the set of vertices reachable from $u$ through still-undiscovered territory at the instant $u$ is opened}. This is why DFS subtrees correspond to ``reachable regions,'' and it is the lever behind Tarjan's strongly-connected-components algorithm (Chapter~5): if you arrange the DFS so that the white-reachable region from $u$ is precisely $u$'s SCC, the subtree hands you the component for free. It is also why a single DFS in a DAG, ordered by finishing time, produces a topological order (Chapter~4).
\end{whybox}

\section{Edge classification}
\label{ch03:edge-types}

Run DFS and every edge of $\Gph$ falls into one of four types, determined by the relationship between the endpoints in the DFS forest. For a directed edge $(u,v)$:

\begin{itemize}
  \item \textbf{Tree edge:} $(u,v)$ was used to first discover $v$; $v$ was white when scanned. These edges form the DFS forest.
  \item \textbf{Back edge:} $v$ is an ancestor of $u$ in the forest ($u$ is in $v$'s subtree). Equivalently $v$ is gray when $(u,v)$ is scanned. A self-loop counts as a back edge.
  \item \textbf{Forward edge:} $v$ is a proper, non-child descendant of $u$ --- reached by an already-built path inside $u$'s subtree. Then $v$ is black when scanned and $d[u]<d[v]$.
  \item \textbf{Cross edge:} $u$ and $v$ are unrelated (neither is an ancestor of the other), and the edge points from a later-finished subtree to an earlier-finished one. Then $v$ is black when scanned and $d[v]<d[u]$.
\end{itemize}

\begin{figure}[H]
\centering
\begin{tikzpicture}[
  >={Stealth[round]}, level distance=14mm, sibling distance=20mm,
  vtx/.style={circle,draw=nodeblue,thick,fill=blue!4,minimum size=8mm,font=\small},
  tree/.style={draw=structurepurple,very thick,->},
  back/.style={draw=trapred,thick,->,bend right=45},
  fwd/.style={draw=quotientorange,thick,->,bend left=55},
  cross/.style={draw=targetgreen,thick,->,dotted}]
  \node[vtx] (a) {a};
  \node[vtx] (b) [below left=of a] {b};
  \node[vtx] (e) [below right=of a] {e};
  \node[vtx] (c) [below=of b] {c};
  \node[vtx] (d) [below=of c] {d};
  \path
    (a) edge[tree] (b)
    (b) edge[tree] (c)
    (c) edge[tree] (d)
    (a) edge[tree] (e)
    (d) edge[back] (b)
    (a) edge[fwd] (c)
    (e) edge[cross] (c);
  \node[anchor=center, font=\scriptsize, text=gray!80!black] at (0, -4.5) {
    \textbf{Classification:}\quad
    \tikz[baseline=-0.5ex]{\draw[structurepurple,very thick,->] (0,0) -- (0.5,0);} Tree \quad
    \tikz[baseline=-0.5ex]{\draw[trapred,thick,->,dashed] (0,0) -- (0.5,0);} Back ($d\!\to\!b$) \quad
    \tikz[baseline=-0.5ex]{\draw[quotientorange,thick,->] (0,0) -- (0.5,0);} Forward ($a\!\to\!c$) \quad
    \tikz[baseline=-0.5ex]{\draw[targetgreen,thick,->,dotted] (0,0) -- (0.5,0);} Cross ($e\!\to\!c$)
  };
\end{tikzpicture}
\caption{The four directed edge types on one forest. Back edges go up to an ancestor; forward edges skip down to a deeper descendant; cross edges jump between unrelated subtrees, always toward something already finished.}
\label{ch03:fig-edgetypes}
\end{figure}

\subsection{Reading the type off the timestamps}

The branch logic in Algorithm~\ref{ch03:alg-dfs} classifies by color at scan time, but it is often more convenient to classify \emph{after the fact} from the final $d/f$ stamps alone. The following lemma makes that precise and, in passing, proves the algorithm's color-based test is correct.

\begin{lemma}[Timestamp signature of edge types]\label{ch03:lem-signature}
Let $(u,v)$ be a directed edge of $\Gph$. After DFS completes:
\begin{enumerate}
  \item[(T/F)] $(u,v)$ is a tree or forward edge $\iff$ $d[u]<d[v]<f[v]<f[u]$ (i.e.\ $I_v\subset I_u$);
  \item[(B)] $(u,v)$ is a back edge $\iff$ $d[v]\le d[u]<f[u]\le f[v]$ (i.e.\ $I_u\subseteq I_v$); a self-loop gives equality;
  \item[(C)] $(u,v)$ is a cross edge $\iff$ $d[v]<f[v]<d[u]<f[u]$ (i.e.\ $I_v$ entirely precedes $I_u$).
\end{enumerate}
\end{lemma}

\begin{proof}
By the parenthesis theorem the two intervals are disjoint or nested, giving exactly three layouts for distinct vertices: $I_v\subset I_u$, $I_u\subset I_v$, or disjoint (and for disjoint, $I_v$ before $I_u$ or after). We match each edge type to a layout.

If $(u,v)$ is a \emph{tree} edge, $v$ is a child of $u$, so by Corollary~\ref{ch03:cor-nest} $I_v\subset I_u$, signature (T/F). If $(u,v)$ is a \emph{forward} edge, $v$ is a deeper descendant of $u$, again $I_v\subset I_u$. Conversely, suppose $I_v\subset I_u$; then $v$ is a descendant of $u$, so when $u$'s list is scanned and reaches $v$, either $v$ is still white (giving a tree edge) or already discovered inside $u$'s subtree (giving a forward edge). Either way it is tree-or-forward. This proves (T/F).

If $(u,v)$ is a \emph{back} edge, $v$ is an ancestor of $u$ (or $u=v$), so $I_u\subseteq I_v$. Conversely if $I_u\subseteq I_v$ with $u\ne v$ then $u$ is a descendant of $v$; the edge $u\to v$ points up to an ancestor, which is exactly a back edge. Self-loops have $I_u=I_v$. This proves (B).

The only remaining layout is $I_v$ and $I_u$ disjoint. If $I_u$ preceded $I_v$ (so $f[u]<d[v]$), then when $u$'s list was scanned --- during $I_u$, before $d[v]$ --- the vertex $v$ was still white, and DFS would have made $v$ a child of $u$, forcing $I_v\subset I_u$, a contradiction. Hence the only consistent disjoint layout for an actual edge is $I_v$ before $I_u$, i.e.\ $d[v]<f[v]<d[u]$. That is the cross signature (C), and it is what is left over.
\end{proof}

\begin{trapbox}
A common mistake is to believe ``edge $(u,v)$ with $d[u]<d[v]$ is a tree or forward edge.'' That is necessary but not sufficient --- the nesting $f[v]<f[u]$ is what seals it. The clean rule is the interval rule of Lemma~\ref{ch03:lem-signature}: compare \emph{both} endpoints' intervals, not just discovery times. Cross edges always satisfy $d[v]<d[u]$, which is why the algorithm's forward/cross tiebreak compares discovery times.
\end{trapbox}

\subsection{Verifying on the trace}

Return to our digraph timestamp intervals:
$1{:}(1,8)$, $2{:}(2,5)$, $3{:}(3,4)$, $4{:}(6,7)$, $5{:}(9,12)$, and $6{:}(10,11)$,
and check Lemma~\ref{ch03:lem-signature} edge by edge.

\begin{center}
\begin{tabular}{@{}llll@{}}
\toprule
edge & $I_u$ & $I_v$ & type \\
\midrule
$1\to2$ & $[1,8]$ & $[2,5]$ & $I_v\subset I_u$, tree \\
$1\to4$ & $[1,8]$ & $[6,7]$ & $I_v\subset I_u$, tree \\
$2\to3$ & $[2,5]$ & $[3,4]$ & $I_v\subset I_u$, tree \\
$3\to1$ & $[3,4]$ & $[1,8]$ & $I_u\subset I_v$, back \\
$4\to3$ & $[6,7]$ & $[3,4]$ & $I_v$ before $I_u$, cross \\
$5\to4$ & $[9,12]$ & $[6,7]$ & $I_v$ before $I_u$, cross \\
$5\to6$ & $[9,12]$ & $[10,11]$ & $I_v\subset I_u$, tree \\
$6\to6$ & $[10,11]$ & $[10,11]$ & $I_u=I_v$, back (self-loop) \\
\bottomrule
\end{tabular}
\end{center}

Every label matches what we recorded during the live trace --- the after-the-fact interval rule and the on-the-fly color rule agree, as the lemma promised. There are no forward edges in this particular graph; add the edge $1\to3$ and it would be classified forward ($I_3=[3,4]\subset I_1=[1,8]$ but $3$ is not a child of $1$).

\section{Cycle detection}
\label{ch03:cycles}

DFS turns cycle detection into a search for one edge type. The directed and undirected cases differ in a way that trips up nearly everyone the first time, so we treat them separately and prove both.

\subsection{Directed graphs: back edge $\iff$ cycle}

\begin{theorem}[Back edges and directed cycles]\label{ch03:thm-dircycle}
A directed graph $\Gph$ contains a (directed) cycle if and only if any DFS of $\Gph$ produces a back edge.
\end{theorem}

\begin{proof}
($\Leftarrow$) Suppose DFS produces a back edge $(u,v)$. Then $v$ is an ancestor of $u$ (or $u=v$), so there is a directed tree path $v\rightsquigarrow u$. Appending the edge $u\to v$ closes a directed cycle $v\rightsquigarrow u\to v$.

($\Rightarrow$) Suppose $\Gph$ has a directed cycle $C$. Let $v$ be the first vertex of $C$ that DFS discovers (the one with smallest $d$ among $C$'s vertices), and let $u\to v$ be the edge of $C$ entering $v$ (its predecessor on the cycle). At time $d[v]$, every other vertex of $C$ is white, and the cycle gives a directed path $v\rightsquigarrow u$ all of whose vertices lie on $C$, hence are white at $d[v]$. By the white-path theorem, $u$ becomes a descendant of $v$. But then the edge $u\to v$ points from a descendant up to its ancestor $v$, which is a back edge.
\end{proof}

This gives a clean linear-time test: run DFS; the moment you scan an edge into a \textsc{gray} vertex, report a cycle (and you can recover it by walking the gray stack from $u$ up to $v$). No need to finish the whole search.

\begin{algorithm}[H]
\DontPrintSemicolon
\SetKwInOut{Input}{input}\SetKwInOut{Output}{output}
\Input{Directed graph $\Gph=(V,E)$.}
\Output{\textsc{true} with a witness cycle if one exists, else \textsc{false}.}
\BlankLine
color all vertices \textsc{white}\;
\ForEach{$u\in V$}{\lIf{$\mathit{color}[u]=\textsc{white}$}{\lIf{\textsc{Probe}$(u)$}{\Return \textsc{true}}}}
\Return \textsc{false}\;
\BlankLine
\SetKwProg{Fn}{procedure}{}{}
\Fn{\textsc{Probe}$(u)$}{
  $\mathit{color}[u]\gets\textsc{gray}$\;
  \ForEach{$v\in\mathit{Adj}[u]$}{
    \uIf{$\mathit{color}[v]=\textsc{gray}$}{\Return \textsc{true} \textnormal{(back edge $u\to v$; cycle is the gray-stack segment $v\!\rightsquigarrow\!u$ plus $u\!\to\!v$)}}
    \uElseIf{$\mathit{color}[v]=\textsc{white}$}{\lIf{\textsc{Probe}$(v)$}{\Return \textsc{true}}}
  }
  $\mathit{color}[u]\gets\textsc{black}$;\quad \Return \textsc{false}\;
}
\caption{Directed cycle detection via gray vertices}
\label{ch03:alg-dircycle}
\end{algorithm}

\subsection{Undirected graphs: the parent trap}

\begin{trapbox}
In an undirected graph, the tree edge $u-v$ you just came down appears again when you scan $v$'s list (because the edge $\{u,v\}$ lives on both lists). If you blindly call ``an edge to a gray vertex'' a back edge, you will declare a cycle for \emph{every} tree edge --- a catastrophic false positive. The fix: when scanning $v$, ignore the single edge back to its DFS parent $\pi[v]=u$. Only an edge to a gray vertex that is \emph{not} the parent signals a real cycle.
\end{trapbox}

There is a sharper subtlety than ``skip the parent'': you must skip the parent \emph{edge}, not every edge to the parent vertex. If $\Gph$ has \emph{parallel} edges (a multigraph) between $u$ and $v$, two distinct edges connect them, and the second one is a genuine cycle of length two even though it lands on the parent vertex. Robust code skips by edge identity (or remembers it has already used the one parent edge), not by vertex name.

\begin{theorem}[Undirected cycles via non-parent back edges]\label{ch03:thm-undircycle}
Let $\Gph$ be a simple undirected graph (no self-loops, no parallel edges). $\Gph$ contains a cycle if and only if a DFS of $\Gph$ encounters an edge $\{u,v\}$ with $v$ gray and $v\ne\pi[u]$.
\end{theorem}

\begin{proof}
First note that in undirected DFS \emph{every} non-tree edge is a back edge: there are no forward or cross edges. Indeed, suppose $\{u,v\}$ is scanned from $u$ and $v$ is not white. If $v$ were black, then $f[v]$ is in the past, so $v$ finished before $u$'s scan; but the edge $\{u,v\}$ is also on $v$'s list, and when $v$ scanned it earlier, $u$ was either white (making $\{u,v\}$ a tree edge with $u$ a child of $v$) or gray (making $u$ an ancestor of $v$) --- in both cases $u$ and $v$ are in an ancestor/descendant relation, so the edge is a tree or back edge, not a cross edge. Hence the only non-tree possibility when scanning from $u$ is $v$ gray: an ancestor of $u$. So undirected DFS yields only tree and back edges (Proposition~\ref{ch03:prop-undir-edges} restates this).

($\Leftarrow$) A back edge $\{u,v\}$ with $v\ne\pi[u]$ and $v$ an ancestor of $u$ means there is a tree path $v\rightsquigarrow u$ of length $\ge 2$ (since $v$ is not the immediate parent), and adding $\{u,v\}$ closes a cycle of length $\ge 3$.

($\Rightarrow$) Suppose $\Gph$ has a cycle $C=w_0-w_1-\cdots-w_{k-1}-w_0$ with $k\ge 3$. Let $w_i$ be the first vertex of $C$ discovered. Its two cycle-neighbors are white at $d[w_i]$, and every vertex of $C$ is reachable from $w_i$ along $C$ through vertices white at $d[w_i]$; by the white-path theorem all of $C$ lies in $w_i$'s subtree. Now walk the cycle starting from $w_i$ along one arc; the discovery times need not increase monotonically forever, so let $\{x,y\}$ be the first cycle edge (in tree-walk order from $w_i$) whose endpoint $y$ was \emph{already discovered} when $x$ scanned it. Such an edge exists because the cycle returns to the already-discovered $w_i$. At the moment $x$ scans $\{x,y\}$, vertex $y$ is gray (it is an ancestor of $x$ on the tree, being discovered earlier and still open), and $y\ne\pi[x]$ because $\{x,y\}$ is a cycle chord rather than the tree edge by which $x$ was reached. The non-parent gray condition fires.
\end{proof}

\begin{algorithm}[H]
\DontPrintSemicolon
\SetKwInOut{Input}{input}\SetKwInOut{Output}{output}
\Input{Undirected graph $\Gph=(V,E)$ as adjacency lists of (neighbor, edge-id).}
\Output{\textsc{true} if $\Gph$ has a cycle, else \textsc{false}.}
\BlankLine
color all \textsc{white}\;
\ForEach{$s\in V$}{\lIf{$\mathit{color}[s]=\textsc{white}$}{\lIf{\textsc{Probe}$(s,\textsc{nil})$}{\Return \textsc{true}}}}
\Return \textsc{false}\;
\BlankLine
\SetKwProg{Fn}{procedure}{}{}
\Fn{\textsc{Probe}$(u,e_{\mathit{in}})$}{
  $\mathit{color}[u]\gets\textsc{gray}$\;
  \ForEach{$(v,e)\in\mathit{Adj}[u]$}{
    \lIf{$e=e_{\mathit{in}}$}{\textbf{continue} \textnormal{(skip the exact edge we entered through)}}
    \uIf{$\mathit{color}[v]=\textsc{gray}$}{\Return \textsc{true}}
    \uElseIf{$\mathit{color}[v]=\textsc{white}$}{\lIf{\textsc{Probe}$(v,e)$}{\Return \textsc{true}}}
  }
  $\mathit{color}[u]\gets\textsc{black}$;\quad \Return \textsc{false}\;
}
\caption{Undirected cycle detection (skip the entering edge by identity)}
\label{ch03:alg-undircycle}
\end{algorithm}

\begin{labbox}
Trace Algorithm~\ref{ch03:alg-undircycle} on the triangle $a-b-c-a$. \textsc{Probe}$(a,\bot)$ grays $a$, recurses on $b$ via edge $ab$. \textsc{Probe}$(b,ab)$ grays $b$, skips $ab$, recurses on $c$ via $bc$. \textsc{Probe}$(c,bc)$ grays $c$, skips $bc$, scans $ca$: neighbor $a$ is gray and edge $ca\ne bc$ --- back edge, return \textsc{true}. Now contrast the path $a-b$: \textsc{Probe}$(a,\bot)$ recurses on $b$ via $ab$; \textsc{Probe}$(b,ab)$ skips $ab$, finds no other neighbor, returns \textsc{false}. No false alarm.
\end{labbox}

\begin{whybox}
Why does the directed test need \emph{no} parent skip while the undirected one does? In a digraph the edge $u\to v$ and the edge $v\to u$ are different objects; coming down a tree edge $\pi[v]\to v$ does not put a $v\to\pi[v]$ edge on $v$'s list. In an undirected graph the single edge $\{u,v\}$ is stored on both endpoints' lists, so traversing it once leaves a phantom ``reverse'' copy that must be filtered out. The asymmetry is entirely an artifact of how undirected edges are represented, and it is the prototype of a dozen later ``don't walk back up the edge you came from'' guards (bridges, biconnectivity, Euler tours).
\end{whybox}

\subsection{Counting independent cycles}

DFS also hands you the dimension of the cycle space cheaply. In an undirected connected graph the DFS tree has $n-1$ edges; every one of the remaining $m-(n-1)$ edges is a back edge, and each back edge closes exactly one fundamental cycle. The number $m-n+1$ is the \emph{cyclomatic number} (first Betti number) of the graph, and the fundamental cycles of the back edges form a basis of the cycle space.

\begin{corollary}\label{ch03:cor-betti}
For an undirected graph with $c$ connected components, the number of independent cycles (back edges across a DFS forest) is $m-n+c$. The graph is a forest iff this is $0$, i.e.\ iff DFS finds no back edge.
\end{corollary}

\begin{proof}
The DFS forest has one tree per component, with $n-c$ tree edges in total (a tree on $n_i$ vertices has $n_i-1$ edges, summed over components). By Proposition~\ref{ch03:prop-undir-edges} every remaining edge is a back edge, and there are $m-(n-c)=m-n+c$ of them. Each is the unique non-tree edge of its fundamental cycle, and standard linear algebra over $\mathrm{GF}(2)$ shows these fundamental cycles are independent and span the cycle space, whose dimension is therefore $m-n+c$.
\end{proof}

\subsection{A worked undirected trace with low-links}
\label{ch03:undir-trace}

The directed example showed all four edge types; an undirected example shows the cleaner two-type world of Proposition~\ref{ch03:prop-undir-edges} and previews the $\mathit{low}$-value bookkeeping that powers Chapter~10. Take the undirected graph on five vertices $\{a,b,c,d,e\}$ with edges $\{a,b\},\{b,c\},\{c,a\},\{c,d\},\{d,e\}$ --- a triangle $a\,b\,c$ with a tail $c-d-e$ hanging off it.

\begin{figure}[H]
\centering
\begin{tikzpicture}[
  node distance=18mm,
  vtx/.style={circle,draw=nodeblue,very thick,fill=blue!4,minimum size=9mm,font=\small\bfseries},
  tree/.style={draw=structurepurple,very thick},
  back/.style={draw=trapred,thick,dashed},
  stamp/.style={font=\scriptsize,text=black!70}]
  \node[vtx,label={[stamp]above:{$d{=}1$}}] (a) {a};
  \node[vtx,label={[stamp]above:{$d{=}2$}}] (b) [right=of a] {b};
  \node[vtx,label={[stamp]below right:{$d{=}3$}}] (c) [below right=12mm of a] {c};
  \node[vtx,label={[stamp]right:{$d{=}4$}}] (d) [below=of c] {d};
  \node[vtx,label={[stamp]right:{$d{=}5$}}] (e) [below=of d] {e};
  \path
    (a) edge[tree] (b)
    (b) edge[tree] (c)
    (c) edge[back] (a)
    (c) edge[tree] (d)
    (d) edge[tree] (e);
\end{tikzpicture}
\caption{Undirected trace started at $a$. Tree edges (purple, solid) $a-b$, $b-c$, $c-d$, $d-e$; the lone non-tree edge $c-a$ (red, dashed) is a back edge to the ancestor $a$. There are no forward or cross edges, as Proposition~\ref{ch03:prop-undir-edges} guarantees.}
\label{ch03:fig-undir}
\end{figure}

Starting the DFS at $a$ and scanning each list alphabetically, the discovery order is $a,b,c,d,e$. The edge $c-a$ is scanned from $c$ while $a$ is still gray and $a\ne\pi[c]=b$: a genuine back edge, so the triangle is a cycle. Now compute $\mathit{low}[v]$ bottom-up, where $\mathit{low}[v]$ is the smallest discovery time reachable from $v$'s subtree using tree edges downward and at most one back edge:

\begin{center}
\begin{tabular}{@{}lllll@{}}
\toprule
$v$ & $d[v]$ & back edge from subtree & $\mathit{low}[v]$ & tree edge $(\pi[v],v)$ a bridge? \\
\midrule
$e$ & $5$ & none & $5$ & $\mathit{low}[e]=5>d[d]=4$ \;\textbf{bridge} \\
$d$ & $4$ & none & $4$ & $\mathit{low}[d]=4>d[c]=3$ \;\textbf{bridge} \\
$c$ & $3$ & $c{-}a$ reaches $d{=}1$ & $1$ & $\mathit{low}[c]=1\not>d[b]=2$ \\
$b$ & $2$ & via $c$, reaches $1$ & $1$ & $\mathit{low}[b]=1\not>d[a]=1$ \\
$a$ & $1$ & --- (root) & $1$ & root, not applicable \\
\bottomrule
\end{tabular}
\end{center}

Read off the bridge criterion $\mathit{low}[v]>d[\pi[v]]$: the tail edges $c-d$ and $d-e$ are bridges (removing either disconnects the graph), while the triangle edges $a-b$, $b-c$, $c-a$ are not (each lies on the cycle $a\,b\,c$). That is the entire content of Chapter~10's bridge-finding algorithm, computed here by hand on one DFS pass --- another structure read straight off the timestamps.

\section{Directed versus undirected: what changes}
\label{ch03:dir-vs-undir}

It is worth pinning down precisely which edge types can occur, because later chapters silently rely on these restrictions.

\begin{proposition}[Undirected DFS has only tree and back edges]\label{ch03:prop-undir-edges}
A DFS of an undirected graph classifies every edge as either a tree edge or a back edge; forward and cross edges never occur.
\end{proposition}

\begin{proof}
Established inside the proof of Theorem~\ref{ch03:thm-undircycle}: an edge to a non-white vertex must reach a gray ancestor, hence a back edge; if it reached white it is a tree edge. The orientation ``forward'' (down to a deeper descendant) is impossible because such an edge would have been used as a tree edge from the shallower endpoint when that endpoint scanned it first; and ``cross'' (between unrelated subtrees) is impossible because an undirected edge between $u$ and $v$ forces one to be discovered while scanning the other, creating an ancestor relationship.
\end{proof}

In a \emph{directed} graph all four types appear. The key structural facts:

\begin{itemize}
  \item \textbf{No back edge} $\iff$ the digraph is acyclic (Theorem~\ref{ch03:thm-dircycle}). A DAG's DFS therefore has only tree, forward, and cross edges, and \emph{ordering vertices by decreasing finish time topologically sorts it} (Chapter~4).
  \item \textbf{Cross edges always go ``right to left''} in time: for a cross edge $(u,v)$, $f[v]<d[u]$, so $v$'s whole subtree finished before $u$ began. This monotonicity is what makes Kosaraju's and Tarjan's SCC algorithms work (Chapter~5).
  \item In undirected graphs, the absence of cross edges is what lets the single number $\mathit{low}[v]$ --- the earliest discovery time reachable from $v$'s subtree via one back edge --- summarize all the connectivity you need for bridges and articulation points (Chapter~10).
\end{itemize}

\begin{historybox}
Depth-first search is ancient as a maze procedure --- Charles Pierre Tr\'emaux described a string-and-mark method for traversing mazes in the 1880s, and \'Edouard Lucas reported it. Its transformation from ``a way to walk a maze'' into ``a way to compute global graph structure in linear time'' is due chiefly to Robert Tarjan and John Hopcroft around 1971--1973. Tarjan's insight was that the discovery/finish timestamps and the $\mathit{low}$-value bookkeeping expose strong connectivity, biconnectivity, planarity, and dominators, all in $\Oh(n+m)$. The parenthesis structure is implicit in any stack-based recursion, but naming it and mining it is what turned DFS into the Swiss-army knife of graph algorithms. Almost every algorithm in Chapters 4, 5, 10, and 21 is a Tarjan-style reading of the timestamps.
\end{historybox}

\section{The structure DFS exposes}
\label{ch03:structure}

We close by collecting the structural payoffs, because the entire reason to prefer DFS over BFS for these problems is that DFS produces \emph{more} than a traversal --- it produces a forest plus a total order on vertices (by $d$ or by $f$) whose interaction encodes deep facts.

\begin{applicationbox}
\textbf{Topological sort (Ch.~4).} Decreasing finish time is a topological order of any DAG. One DFS, then read $f[\cdot]$ in reverse. \\[2pt]
\textbf{Strongly connected components (Ch.~5).} Kosaraju: DFS for finish order, then DFS the transpose in that order; each tree is an SCC. Tarjan: a single DFS with a $\mathit{low}$-link and an auxiliary stack pops one SCC at a time. Both are timestamp readings. \\[2pt]
\textbf{Bridges and articulation points (Ch.~10).} $\mathit{low}[v]=\min$ of $d[v]$ and the discovery times reachable by one back edge from $v$'s subtree. An edge $(u,v)$ is a bridge iff $\mathit{low}[v]>d[u]$. \\[2pt]
\textbf{Euler tour / subtree-as-interval (Ch.~21,~23).} The pair $(d[v],f[v])$ is exactly an interval; ``$v$ in $u$'s subtree'' becomes ``$I_v\subset I_u$,'' turning subtree queries into range queries on a segment tree. \\[2pt]
\textbf{Dominator trees, planarity testing.} Both are built on DFS numbering and the low-link/semidominator machinery.
\end{applicationbox}

\begin{connectionbox}
The interval picture connects forward to the linear-algebra companion. Order vertices by discovery time and the back edges of an undirected DFS are exactly the entries that make the graph's structure differ from that of a tree; the tree edges form a spanning tree whose incidence matrix has full rank $n-1$, and the $m-n+1$ back edges index a basis of the cycle space --- the kernel of the boundary map, dual to the cut space spanned by the tree. When we reach the graph Laplacian $\Lap$ (Ch.~31) and electrical flows (Ch.~35,~41), the cycle/cut decomposition that DFS exposes here reappears as the orthogonal decomposition of flows into circulations and potential-driven currents. The humble back edge is the first appearance of a homology class.
\end{connectionbox}

\begin{recoverybox}
The mental model to keep: DFS does not merely visit vertices, it \emph{nests} them. Two numbers per vertex, laminar by the parenthesis theorem, and four edge types determined by how the intervals sit relative to each other. Cycle detection is ``find a back edge'' (with a parent guard when edges are undirected). Everything heavier --- SCCs, bridges, topological order, Euler tours --- is a disciplined reading of those same two numbers.
\end{recoverybox}

\section{Exercises}
\label{ch03:exercises}

\begin{exercise}[Trace and classify]\label{ch03:ex-trace}
Run DFS on the digraph with edges $a\to b,\ a\to c,\ b\to d,\ c\to d,\ d\to a,\ c\to e,\ e\to c$, visiting roots and adjacency lists in alphabetical order. Record $d[\cdot]$ and $f[\cdot]$, draw the DFS forest, and label every edge as tree/back/forward/cross. Verify each label with the interval rule of Lemma~\ref{ch03:lem-signature}.
\end{exercise}

\begin{exercise}[Parenthesis string]\label{ch03:ex-paren}
Prove that the sequence of discovery and finish events, written as labeled parentheses, is a balanced parenthesization, and that conversely any balanced parenthesization on $n$ symbols is realizable as the event sequence of some DFS on some graph. How many distinct DFS forest shapes (nesting structures) on $n$ labeled vertices are there?
\end{exercise}

\begin{exercise}[Correct iterative DFS]\label{ch03:ex-iterative}
Write an explicit-stack DFS that reproduces $d[\cdot]$, $f[\cdot]$, $\pi[\cdot]$, and the edge classification of the recursive version \emph{exactly}, including correct finish times. The trap: each stack frame must store an iterator into its adjacency list so it resumes where it left off, and the finish timestamp must be set when the frame is popped. Test it against Algorithm~\ref{ch03:alg-dfs} on the graph of Figure~\ref{ch03:fig-digraph}.
\end{exercise}

\begin{exercise}[Break the undirected detector]\label{ch03:ex-break}
Exhibit an undirected \emph{multigraph} (two vertices, two parallel edges) on which the ``skip the parent vertex'' version of cycle detection wrongly reports no cycle, but the ``skip the parent edge by identity'' version of Algorithm~\ref{ch03:alg-undircycle} correctly reports one. Explain precisely which line saves it.
\end{exercise}

\begin{exercise}[Forward edges require direction]\label{ch03:ex-fwd}
Prove that a forward edge $(u,v)$ can exist only in a directed graph (cf.\ Proposition~\ref{ch03:prop-undir-edges}), and give the smallest directed graph (fewest edges) whose DFS, under a stated visiting order, produces a forward edge. Confirm with timestamps.
\end{exercise}

\begin{exercise}[Cross-edge monotonicity]\label{ch03:ex-cross}
Prove that for every cross edge $(u,v)$ in a directed DFS we have $f[v]<d[u]$, and hence $f[v]<f[u]$. Then show that if a digraph's DFS has no back edge, ordering vertices by decreasing $f$ gives an ordering in which every edge goes from earlier to later --- a topological sort. (This is the seed of Chapter~4.)
\end{exercise}

\begin{exercise}[Counting components and cycles]\label{ch03:ex-betti}
Using only the DFS forest, give an $\Oh(n+m)$ algorithm that returns the number of connected components $c$ of an undirected graph and the dimension $m-n+c$ of its cycle space, and outputs an explicit fundamental cycle for one chosen back edge. Prove correctness via Corollary~\ref{ch03:cor-betti}.
\end{exercise}

\begin{exercise}[Order-dependence of the forest]\label{ch03:ex-order}
The DFS forest and the edge classification depend on the visiting order, but some facts are invariant. Prove that whether a directed graph has a back edge (equivalently, whether it has a cycle) is independent of the visiting order, even though \emph{which} edges are back edges is not. Give a 3-vertex digraph where two visiting orders turn the same edge from a cross edge into a forward edge.
\end{exercise}

\begin{exercise}[White-path in action]\label{ch03:ex-whitepath}
Use the white-path theorem (Theorem~\ref{ch03:thm-whitepath}) to prove the following without re-running DFS: if $u$ is discovered before $v$ and there is a directed path $u\rightsquigarrow v$ whose internal vertices are all undiscovered at time $d[u]$, then $f[v]<f[u]$. Where does the argument break if some internal vertex is already gray at $d[u]$?
\end{exercise}

\begin{exercise}[Low-link preview]\label{ch03:ex-low}
For an undirected graph define $\mathit{low}[v]=\min\big(\{d[v]\}\cup\{d[w] : \{x,w\}\text{ a back edge with }x\text{ in }v\text{'s subtree}\}\big)$. Give a single-pass DFS that computes $\mathit{low}[v]$ for all $v$ in $\Oh(n+m)$, and prove that a tree edge $(u,v)$ is a bridge iff $\mathit{low}[v]>d[u]$. (Full treatment in Chapter~10; here just establish the recurrence and the bridge criterion.)
\end{exercise}

\begin{exercise}[Undirected bridges by hand]\label{ch03:ex-bridges}
Add the edge $\{e,a\}$ to the graph of Figure~\ref{ch03:fig-undir}, turning the tail into a second cycle $a\,c\,d\,e$. Recompute every $\mathit{low}[v]$ from scratch and show that \emph{no} tree edge now satisfies $\mathit{low}[v]>d[\pi[v]]$, so the augmented graph has no bridges. Explain in one sentence, using the $\mathit{low}$ values, why a graph with no bridges cannot be disconnected by deleting any single edge.
\end{exercise}

\begin{exercise}[Cursor correctness]\label{ch03:ex-cursor}
Prove that Algorithm~\ref{ch03:alg-iterative} assigns exactly the same $d[\cdot]$ and $f[\cdot]$ as the recursive Algorithm~\ref{ch03:alg-dfs} when both start from the same root and scan adjacency lists in the same order. The crux is the \textbf{else} branch: argue that a frame is popped (and its $f$ stamp set) at precisely the clock tick when the matching recursive call would return. Then modify the skeleton to also emit the four edge labels, and confirm it reproduces the table of Section~\ref{ch03:edge-types}.
\end{exercise}

\begin{exercise}[Deep recursion stress test]\label{ch03:ex-deep}
Construct the family of graphs on which recursive DFS recurses to depth exactly $n$, and a family on which it recurses to depth $\Oh(\log n)$ for all visiting orders. For the first family, estimate the largest $n$ your language's default stack survives, then confirm your iterative DFS from Exercise~\ref{ch03:ex-iterative} handles $10\times$ that.
\end{exercise}
